\documentclass{beamer}
\usepackage[spanish]{babel}
\usepackage{graphicx}
\usepackage{booktabs}
\usepackage{amsmath}

\title{Optimización de Portafolio de Inversión}
\subtitle{Solución para OptimaBattle Arena}
\author{Equipo de Optimización}
\date{\today}

\begin{document}

\frame{\titlepage}

\begin{frame}
\frametitle{Introducción}
\begin{itemize}
\item Implementación de un optimizador de portafolio para el torneo OptimaBattle
\item Objetivo: Maximizar retorno ajustado al riesgo bajo restricciones específicas
\item Método: Programación convexa con variables enteras (MIQP)
\item Lenguaje: Python con librerías especializadas (CVXPY, Pandas, NumPy)
\end{itemize}
\end{frame}

\begin{frame}
\frametitle{Método de Optimización}
\begin{block}{Problema de Optimización Mixta-Entera}
\begin{itemize}
\item \textbf{Tipo}: Programación Cuadrática Entera Mixta (MIQP)
\item \textbf{Solver}: ECOS\_BB (Branch-and-Bound para problemas convexos)
\item \textbf{Variables}:
\begin{itemize}
\item $w_i$: Variables continuas (pesos del portafolio)
\item $y_i$: Variables binarias (selección de activos)
\end{itemize}
\end{itemize}
\end{block}

\begin{equation*}
\text{Maximizar } U = \sum_{i=1}^n r_i w_i - \lambda \sum_{i=1}^n \sigma_i^2 w_i^2
\end{equation*}
\end{frame}

\begin{frame}
\frametitle{Restricciones Implementadas}
\begin{table}[]
\small
\begin{tabular}{ll}
\toprule
\textbf{Restricción} & \textbf{Implementación} \\
\midrule
Presupuesto total & $\sum w_i = 1$ \\
No ventas en corto & $w_i \geq 0$ \\
Diversificación sectorial & $\sum_{i \in S_j} w_i \leq 0.3$ \\
Mínimo de activos & $\sum y_i \geq 5$ \\
Límite de riesgo sistemático & $\sum \beta_i w_i \leq 1.2$ \\
Inversión mínima por activo & $w_i \cdot B \geq m_i \cdot y_i$ \\
\bottomrule
\end{tabular}
\end{table}
\end{frame}

\begin{frame}
\frametitle{Flujo de Solución}
\begin{enumerate}
\item \textbf{Verificación de datos}: Chequeo de archivo y validación de rangos
\item \textbf{Preparación de parámetros}: Normalización de retornos y volatilidades
\item \textbf{Formulación del problema}: Construcción de función objetivo y restricciones
\item \textbf{Solución}: Uso de solver ECOS\_BB
\item \textbf{Post-procesamiento}: Ajuste a números enteros de acciones
\item \textbf{Cálculo de métricas}: Retorno, riesgo, beta y puntaje
\end{enumerate}
\end{frame}

\begin{frame}
\frametitle{Resultados Obtenidos}
\begin{columns}
\begin{column}{0.5\textwidth}
\begin{block}{Métricas Clave}
\begin{itemize}
\item Retorno esperado: 12.86\%
\item Volatilidad: 88.45\%
\item Beta: 1.20
\item Puntaje: -376.35
\item Activos seleccionados: 29
\end{itemize}
\end{block}
\end{column}
\begin{column}{0.5\textwidth}
\begin{block}{Distribución Sectorial}
\begin{itemize}
\item Sector 1: 29.9\%
\item Sector 2: 4.3\%
\item Sector 3: 30.0\%
\item Sector 4: 5.7\%
\item Sector 5: 30.0\%
\end{itemize}
\end{block}
\end{column}
\end{columns}
\end{frame}

\begin{frame}
\frametitle{Análisis de Resultados}
\begin{itemize}
\item \textbf{Cumplimiento de restricciones}:
\begin{itemize}
\item Todas las restricciones fueron satisfechas (beta exactamente en 1.2)
\item Distribución sectorial cerca del límite máximo (30\%)
\end{itemize}
\item \textbf{Problemas identificados}:
\begin{itemize}
\item Puntaje negativo debido a alta volatilidad
\item Posible sobre-diversificación (29 activos)
\end{itemize}
\item \textbf{Mejoras potenciales}:
\begin{itemize}
\item Ajustar parámetro de aversión al riesgo ($\lambda$)
\item Considerar liquidez en la función objetivo
\item Limitar número máximo de activos
\end{itemize}
\end{itemize}
\end{frame}

\begin{frame}
\frametitle{Conclusiones}
\begin{itemize}
\item Implementación exitosa del modelo de optimización
\item Solución cumple con todas las restricciones del problema
\item Alta volatilidad sugiere necesidad de ajustar parámetros
\item El enfoque MIQP permite manejar restricciones complejas
\item Código estructurado y modular facilita modificaciones
\end{itemize}

\begin{block}{Próximos pasos}
\begin{itemize}
\item Pruebas con diferentes parámetros de aversión al riesgo
\item Incorporar restricción de liquidez mínima
\item Optimizar tiempo de ejecución para competencia
\end{itemize}
\end{block}
\end{frame}

\end{document}